\section{Technical Appendix}

\subsection{ResNet-18 Architecture and Feature Fusion}

The classification backbone is ResNet-18 (~11.7M parameters), chosen for its balance between learning capacity and computational efficiency on the Nvidia Jetson Orin Nano. The model was initialised with ImageNet pre-trained weights to speed up convergence.

\subsubsection{Architecture Modifications}

The original architecture was adapted as follows:

\begin{enumerate}
    \item \textbf{Input Layer}: The first convolutional layer was modified from 3-channel RGB to single-channel int8 spectrograms produced via Short-Time Fourier Transform (STFT).
    \item \textbf{Feature Extraction}: The ResNet backbone outputs 512-dimensional feature vectors.
    \item \textbf{Feature Fusion}: 16 hand-crafted statistical features (mean, median, std, RMS, percentiles, skewness, kurtosis, entropy) are batch-normalised and concatenated with the ResNet features, giving a 528-dimensional vector.
    \item \textbf{Classification Head}: $528 \rightarrow 256$ (ReLU, Dropout 0.5) $\rightarrow 16$ classes. An energy score (log-sum-exp of logits) is computed per sample for open-set detection.
\end{enumerate}

\subsubsection{Dataset and Training}

The dataset contains 44,428 samples across 16 classes: CLEAN + 5 jammer types $\times$ 3 power levels (LOW, MID, HIGH) based on Jammer-to-Signal Ratio. Data was split 50/15/35\% (train/val/test, stratified). Training used Adam (lr=0.001), cross-entropy loss, batch size 128, and early stopping (patience=5). Convergence typically occurred within 20 epochs on GPU.

\subsection{In-Set Classification and Open-Set Recognition}

\subsubsection{Closed-Set Baseline}

Without open-set recognition, the model always forces inputs into one of the 16 known classes via softmax. On known samples it performs well, reaching \textbf{93.8\% accuracy} (3712/3958). However, it completely fails on unknown jammer types: all 789 unknown samples in the extended validation set are assigned to a known class with over 99\% average softmax confidence.

The failure is not random — it is highly systematic. \textbf{99.4\% of all unknowns} (784/789) are misclassified as \textbf{TICK\_LOW}, regardless of the actual jammer type. This reveals that unknown signal embeddings consistently fall close to the TICK\_LOW centroid in the 528-dimensional feature space. On the full mixed set of 4747 samples, overall accuracy appears at 78.2\%, but this figure is misleading since unknowns are not distributed randomly — they almost all land on the same wrong class.

This is a critical problem for deployment: encountering a novel jammer produces a confident false positive as TICK\_LOW rather than an alarm. Softmax scores are not suitable for unknown detection, but the underlying embedding space does preserve enough geometric structure for distance-based methods to work — which is what OSR exploits.

\subsubsection{Open-Set Recognition Framework}

OSR operates on the 528-dimensional penultimate feature space and enables both unknown detection and error correction via distance-based decision boundaries.

\textbf{Implementation:}
\begin{enumerate}
    \item \textbf{Centroid Computation}: Class centroids are computed by averaging training embeddings per class.
    \item \textbf{Distance Metric}: L2 distance from each sample's embedding to all class centroids is computed at inference time.
    \item \textbf{ROC-based Threshold Calibration}: A per-class threshold is set by capping the False Positive Rate (unknowns incorrectly accepted) at $\leq 5\%$, while maximising the True Positive Rate (known samples correctly accepted).
    \item \textbf{Three-State Decision Logic}:
    \begin{itemize}
        \item[\textbf{ACCEPT}]: Distance to predicted class $\leq$ threshold $\rightarrow$ keep prediction
        \item[\textbf{CORRECT}]: Distance exceeds threshold but a better class is within its threshold $\rightarrow$ reassign
        \item[\textbf{REJECT}]: No class within threshold $\rightarrow$ classify as UNKNOWN
    \end{itemize}
\end{enumerate}

\subsubsection{Experimental Results}

\begin{table}[h]
\centering
\begin{tabular}{|l|c|}
\hline
\textbf{Metric} & \textbf{Value} \\
\hline
Overall Accuracy & 93.2\% \\
Known Sample Accuracy & 93.8\% \\
Unknown Detection Rate & 90.1\% \\
Samples Within Threshold & 82.7\% \\
Mean Inference Time & 45.80 ms \\
\hline
\end{tabular}
\caption{OSR Performance on 4747 Samples (3958 known + 789 unknown)}
\end{table}

OSR achieves \textbf{90.1\% unknown detection} (711/789), though performance varies by jammer type: LINEAR-WIDE-20, FH-20, and MULTITONE-NARROW-40-20 all reach \textbf{100\% detection}, while HOOKED-SAWTOOTH-3-20 reaches only \textbf{61.6\%} (125/203) due to feature similarity with TICK\_MID. The 9.9\% miss rate (78 samples) reflects cases where unknown embeddings fall geometrically close to TICK\_MID centroids.

Error correction was triggered in 1.8\% of samples (87/4747), with only 9.2\% of those leading to a successful reassignment. This conservative behaviour is intentional — aggressive corrections introduce more errors than they fix. Known sample accuracy is maintained at 93.8\%, confirming that OSR adds unknown detection without degrading performance on known classes. Additionally, 82.7\% of samples fall within their predicted class threshold, and of the 17.3\% that do not, 89.4\% are correctly flagged as UNKNOWN — validating the ROC calibration strategy.

\subsubsection{Key Findings}

OSR successfully addresses the closed-set limitation, achieving strong unknown detection while preserving known-class accuracy. The main remaining weakness is the 61.6\% detection rate on HOOKED-SAWTOOTH-3-20, caused by its embedding proximity to the TICK family — a limitation of the current feature representation rather than of OSR itself. Improving feature extraction or adding more diverse TICK-like training data are natural directions for future work.